\begin{tikzpicture}[scale=1.62]
  % 坐标轴与边界圆
  \draw[axisline] (-0.45,0) -- (3.30,0) node[right] {$x$};
  \draw[axisline] (0,-0.45) -- (0,3.02) node[above] {$y$};
  \draw[aux] (0,0) circle (2.55);
  \node[below left,font=\small] at (0,0) {$O$};
  \node[font=\small,anchor=west] at (2.60,0.34) {边界 $r=a$};

  % 极角与半径
  \draw[thick] (0,0) -- (30:2.55);
  \draw[thick] (0,0) -- (60:2.55);
  \draw (1.00,0) arc (0:30:1.00);
  \node[font=\small] at (15:1.26) {$\varphi$};
  \node[font=\small,rotate=30] at (0.52,0.38) {$r$};

  % 面积元：径向宽 dr，弧长 r dphi
  \fill[gray!25] (30:1.55) arc (30:60:1.55) -- (60:2.15) arc (60:30:2.15) -- cycle;
  \draw[thick] (30:1.55) arc (30:60:1.55);
  \draw[thick] (30:2.15) arc (30:60:2.15);
  \draw[thick] (30:1.55) -- (30:2.15);
  \draw[thick] (60:1.55) -- (60:2.15);

  % 面积元的标注：引线指向灰色小片
  \node[font=\small,anchor=west] (dA) at (-2.00,1.05)
    {$\mathrm{d}A=r\,\mathrm{d}r\,\mathrm{d}\varphi$};
  \draw[gray,thin,->] (dA.east) -- (1.42,1.50);

  % 点 P：膜上的一点
  \coordinate (P) at (60:2.44);
  \fill (P) circle (1.6pt) node[right=3pt,font=\small] {$P$};
\end{tikzpicture}
